\subsection*{2.4. 孤立奇点}

\subsubsection*{2.4.1.}
本节为简便起见, 始终假设 $S$ 为严格 Hensel 离散赋值环谱. 设 $\Lambda$ 为挠性与 $S$ 的剩余特征互素的 Noether 挠环, $f:X\longrightarrow S$ 为有限型态射, $F$ 为 $X$ 上的可构造 $\Lambda$-模层. 若 $f$ 在 $x$ 处光滑, 且 $F$ 在 $x$ 的一个邻域上局部常值, 则称点 $x\in X$ 为 $(X,F)$ 的光滑点. 以 $\Sigma$ 表示 $(X,F)$ 不光滑的点集, 以 $\Sigma_s$ 表示其特殊纤维. 我们知道, 在 $\Sigma_s$ 之外有 $R\Phi(F)=0$.

\subsubsection*{定理 2.4.2.}
\underline{在 $\Sigma_s$ 的每个孤立点 $x$ 处, 模 $R^i\Phi_{\bar\eta}(F)_x$ 均为有限型 $\Lambda$-模, 且在离散赋值环谱变换下保持不变（通过 2.1.7.1）.}

在等特征零的情形中, 这些断言已在 2.1.12 与 2.3.1 中证明, 不必假设 $x$ 孤立.

设 $x$ 为 $\Sigma_s$ 的孤立点. 定理在 $x$ 的邻域中是局部的, 因而可假设 $f$ 固有. 此时有长正合列 (2.1.8.9)
\[
\cdots\longrightarrow H^i(X_s,F)
\longrightarrow H^i(X_{\bar\eta},F)
\longrightarrow
\mathbb H^i\bigl(X_s,R\Phi_{\bar\eta}(F)\bigr)
\longrightarrow\cdots .
\]
在此列中, $H^i(X_s,F)$ 与 $H^i(X_{\bar\eta},F)$ 在离散赋值环谱变换下保持不变, 且均为有限型 $\Lambda$-模（SGA~4, XIV, 1.1）. 因此模 $\mathbb H^i(X_s,R\Phi_{\bar\eta}(F))$ 也有相同性质. 回顾下述引理.

\subsubsection*{引理 2.4.3.}
\underline{设 $(T,\Lambda)$ 为赋环拓扑斯, $Z$ 为 $T$ 的闭子拓扑斯, $i:Z\hookrightarrow T$ 为包含态射.} \underline{则函子}
\[
i_*:D(Z,\Lambda)\longrightarrow D(T,\Lambda)
\]
\underline{给出从 $D(Z,\Lambda)$ 到 $D(T,\Lambda)$ 中由满足 $\mathcal H^i(K)$ 的支集包含于 $Z$ 的复形 $K$ 所成的满子范畴的等价.}

事实上, 对这样的复形 $K$, 态射 $K\longrightarrow i_*i^*K$ 为拟同构.

因此可将 $R\Phi_{\bar\eta}(F)$ 看作
\[
D^b(\Sigma_s,\Lambda)
=D^b(\Sigma_s-\{x\},\Lambda)\oplus D^b(\{x\},\Lambda).
\]
的对象. 取一个直和项, 可知
\[
\mathbb H^i\bigl(\{x\},R\Phi_{\bar\eta}(F)\bigr)
=R^i\Phi_{\bar\eta}(F)_x
\]
具有 2.4.2 中所述的性质. 证明完毕.

\subsubsection*{注 2.4.4.}
同样的证明说明, 若 $K$ 在 $D^+(X,\Lambda)$ 中可构造, 且 $x$ 在 $R\Phi_{\bar\eta}(K)$ 的支集中孤立, 则茎 $R\Phi_{\bar\eta}(K)_x$ 的各上同调模在 $\Lambda$ 上有限型. 若作离散赋值环谱变换后, $x$ 仍为该支集的孤立点, 则这些模在此变换下保持不变.
